\documentclass[10pt]{uhhlet}
\usepackage{xcolor,mhchem,mathpazo,ragged2e}
\usepackage[utf8]{inputenc}
\usepackage[T1]{fontenc}
\usepackage{txfonts}
\newcommand{\alert}[1]{\textcolor{red}{#1}}
\usepackage[UKenglish]{datetime}

% Formal environment
\usepackage[strict]{changepage}
\definecolor{darkblue}{rgb}{0.0, 0.0, 0.55}
\definecolor{formalshade}{rgb}{0.95,0.95,1}
\usepackage{framed}
\newcommand{\ajwt}[1]{{\color{red}{#1}}}
\newenvironment{formal}{%
  \def\FrameCommand{%
    {\color{darkblue}\vrule width 2pt}%
    {\color{formalshade}\vrule width 4pt}%
    \colorbox{formalshade}%
  }%
  \MakeFramed{\advance\hsize-\width\FrameRestore}%
  \noindent\hspace{-4.55pt}% disable indenting first paragraph
  \begin{adjustwidth}{}{7pt}%
  \vspace{-4pt}%
}
{%
  \vspace{4pt}\end{adjustwidth}\endMakeFramed%
  \vspace{-8pt}
}



% Paper specifics
\newcommand{\journal}{Journal}
\newcommand{\papertitle}{Title}

\begin{document}

\newdateformat{UKvardate}{%
\THEDAY\ \monthname[\THEMONTH], \THEYEAR}
\UKvardate\renewcommand{\headrulewidth}{0pt}

\fromwhom{}{Dr.~Johannes T\"olle}{}{Insitut f\"ur Anorganische und Angewandte Chemie}{johannes.toelle@chemie.uni-hamburg.de}{+49 157 79677990}{}
\setsalutation{To}
\renewcommand{\yoursdetails}{Yours sincerely, the authors}

\begin{letter}%
{the Editors of \textit{\journal{}}}

%\opening{Dear Editors,}

\vspace{-1em}
\justifying
Please find attached a manuscript entitled 
\vspace{-1.25em}
\begin{center}
\textbf{\papertitle{}}
\end{center}
\vspace{-1.25em}
that we recently submitted as a \underline{Regular Article} to the \textit{\journal}. 
We thank both reviewers for their constructive comments and for supporting publication of the present manuscript. 
For convenience, changes are highlighted in red in the revised version of the manuscript. We look forward to hearing from you.

\end{letter}
%%%%%%%%%%%%%%%%%%%%%%%%%%%%%%
\clearpage
{\large \bf Response to Reviewer 1:}
%%%%%%%%%%%%%%%%%%%%%%%%%%%%%%
\begin{itemize}  
  \item[] \textit{I believe that the title of the manuscript and/or the contents of the abstract should be modified to reflect that this work deals with the connection between a standard linear coupling electron-boson Hamiltonian that gives a first-order diagrammatic self-energy which is the $G_0W_0$ self-energy and the EOM-ECC on this Hamiltonian. This is simply because I feel that, in its current form, the manuscript is slightly over-reaching in its claims of a fundamental connection between these methods. This is because the crucial contribution of Hedin and the $GW$ approximation was to establish the emergence of a screened interaction starting from the unscreened many-body electronic structure Hamiltonian. I feel that by working with a Hamiltonian that already contains this screening in the form of the effective linear coupling electron-boson Hamiltonian is slightly contrived. In addition, the authors seem unaware that this Hamiltonian has its own diagrammatic self-energy expansion that might agree at first-order with the $G_0W_0$ approximation but is fundamentally different to that of the electronic structure self-energy. I recommend that they consult the literature on the electron-phonon interaction as evidence of the different diagrammatic expansion implied by Eqs 24. They should be aware that inclusion of `vertex' corrections based on this approximation is in fact not systematic will not rigorously converge to the exact spectrum of the electronic structure Hamiltonian. }

  \item[] We respectfully disagree with several points raised by the reviewer. In our view, the manuscript does establish a fundamental connection between the $GW$ approximation and EOM-ECC. To this end, we introduced a set of approximations to the general structure of the underlying Hamiltonian intrinsic to the $GW$ framework (\textit{vide supra}). 
  Within these approximations, we demonstrate explicitly how screening emerges in the EOM-ECC formalism (Sec.~V). Without this Ansatz, the screening of the linear coupling Hamiltonian (Eq.~17) is restricted to a subset of bubble diagrams, corresponding to the TDA, as discussed in Sec.~III. 
  In light of the reviewer's emphasis on Hedin's key contribution (the emergence of the screened interaction), we stress that our work reproduces precisely this feature from an ECC perspective.
  Importantly, the ECC framework is essential in this context, as conventional CC theory cannot account for the full set of time orderings required to describe Hedin's screened interaction (cf. Refs.~104 and 124, and the discussion of the introduction).
\end{itemize}


%%%%%%%%%%%%%%%%%%%%%%%%%%%%%%
{\large \bf Response to Reviewer 2:}
%%%%%%%%%%%%%%%%%%%%%%%%%%%%%%
\begin{itemize}
  \item \textit{In the introduction after Eq. 3, the authors write that "Systematic improvements beyond GW can then be obtained by including vertex corrections in either $P$ (internal corrections) or $\Sigma_{xc}$ (external corrections)": I largely agree with this statement. However, it is worth pointing out that systematic improvements will need vertex corrections both in $P$ and $\Sigma$, otherwise one calculates diagrammatically incomplete expansions. For instance, there seems to be compelling evidence that vertex corrections in $\Sigma$ alone worsen the results compared to GW. I would therefore suggest to rephrase the sentence in a way that reflects this.}
  
  \item[] Thank you for pointing out this potential confusion.
  We have changed the sentence to: ``Systematic improvements beyond $GW$ can then, in principle, be obtained by including vertex corrections in either $P$ (internal corrections) and/or $\Sigma_\text{xc}$ (external corrections) in combination with careful diagrammatic analysis to avoid incomplete diagrammatic expansions.''
\end{itemize}

{\large \bf Additional changes:}

\begin{itemize}
  \item[] \textbf{Whatever has been changed additionally.}
\end{itemize}

\end{document}
